% =============================================================================
% Abstract
% =============================================================================
\chapter*{Abstract}
\addcontentsline{toc}{chapter}{Abstract}

Financial statements summarize firms' past performance, yet corporate filings
contain extensive narrative disclosures written by management. Due to the
complexity and length of these filings, investors may face substantial
information processing costs when evaluating the information embedded in
narrative disclosures. This study investigates whether changes in the language
of corporate filings reveal information about future stock return volatility
that is not captured by traditional financial variables. Using a panel of
approximately 600 large U.S.\ non-financial firms over 2010--2025, I construct
a textual-financial divergence measure that captures the gap between the
direction of narrative change in 10-K filings and the direction of concurrent
financial change. I find that when managerial language diverges from the
underlying fundamentals, future stock return volatility increases, consistent
with the hypothesis that such divergence creates information uncertainty among
investors. The divergence measure predicts volatility beyond both financial
variables and textual features alone, suggesting that the interaction between
words and numbers---not either in isolation---is a distinct source of
information about firm-level risk.

\vspace{1cm}
\noindent\textbf{Keywords:} textual analysis, 10-K filings, SEC EDGAR,
volatility, information uncertainty, narrative divergence, sentiment,
panel data, corporate disclosures
